\chapter*{\chapterstyle{Algèbre II {:} Applications Linéaires}}
Soit deux \(\K\)-espaces vectoriels \(E\) et \(F\), et \(f : E \longrightarrow F \).\<

On dit que \(f\) est une \textbf{application linéaire} ou \textbf{morphisme d'espaces vectoriels} si et seulement si pour tout couple de vecteurs \(u, v \in E\) et pour tout scalaire \(\lambda\) elle vérifie: 

\[
   \mathbox{
      \begin{aligned}
          &\bullet \textbf{ Additivité : } f(u + v) = f(u) + f(v)\\
          &\bullet \textbf{ Homogénéité : } f(\lambda u) = \lambda f(u)\\
      \end{aligned}
   }
\]
On note alors \(\mathcal{L}(E, F)\) l'ensemble des applications linéaires de \(E\) vers \(F\).\<

On appelle \textbf{isomorphisme} tout application linéaire bijective.\+
On appelle \textbf{endomorphisme} tout application linéaire d'un ensemble dans lui-même.\+
On appelle \textbf{automorphisme} tout application linéaire bijective d'un ensemble dans lui-même\footnote[1]{On appelle l'ensemble des automorphismes \textbf{le groupe linéaire} de \(E\) qu'on note \(GL(E)\).}.\<

On appelle \textbf{forme linéaire} tout application linéaire qui a pour ensemble d'arrivée \textbf{le corps de base}\footnote[2]{On appelle l'ensemble des formes linéaires \textbf{l'espace dual} de \(E\) qu'on note \(E^*\).}.

\subsection*{\subsecstyle{Propriétés{:}}}
On s'intéresse au propriétés de l'ensemble \(\mathcal{L}(E, F)\), tout d'abord on peut prouver aisément que c'est \textbf{un sous-espace vectoriel} de l'espace des fonctions, ie il est stable par somme et par multiplication exeterne.\+
En outre, \textbf{la composée de deux applications linéaires est linéaire}.\<

Considérons une application linéaire bijective de \(\mathcal{L}(E, F)\), alors:
\begin{center}
   \textbf{Sa réciproque est nécessairement linéaire\footnote[3]{Découle directement des définitions.}}.
\end{center}

\subsection*{\subsecstyle{Image et Noyau{:}}}
On appelle \textbf{image} d'une application linéaire l'image telle que définie pour les applications classiques, ie:
\[
      \text{Im}(f) := \Bigl\{ f(u) \; ; \; u \in E \Bigl\}
\]
Aussi si \(E = \text{Vect}(e_1, e_2, \ldots, e_n)\) pour une certaine famille \((e_i)_{i\in \N}\), alors on a\footnote[5]{Trivialement à partir de la définition du vect et de l'image.}:
\[
      \mathbox{\text{Im}(f) = f(\text{Vect}(e_1, e_2, \ldots, e_n)) = \text{Vect}(f(e_1), f(e_2), \ldots, f(e_n))}
\]
En particulier, on voit que \(\text{Im}(f)\) est un \textbf{sous-espace vectoriel} de l'ensemble d'arrivée.\<

On appelle \textbf{noyau} d'une application linéaire tout \textbf{les antécédent du vecteur nul}, ie:
\[
      \text{Ker}(f) := \Bigl\{ u \in E \; ; \; f(u) = 0_F \Bigl\} = f^{-1}(\{0_F\})
\]
En particulier, on voit directement que \(\text{Ker}(f)\) est un \textbf{sous-espace vectoriel} de l'ensemble de départ car si l'image de deux vecteurs est nulle l'image de leur somme est nulle par linéarité de \(f\).\<

Une propriété fondamentale du noyau\footnote[6]{De manière générale, cette propriété est vraie pour le noyau de tout morphisme de groupe.} est la suivante:
\begin{center}
   \textbox{Une application linéaire est injective si et seulement si son noyau est réduit au vecteur nul.}   
\end{center}
Intuitivement, on peut alors voir la "taille" du noyau comme \textbf{une mesure de la non-injectivité} d'une application.
\pagebreak

\subsection*{\subsecstyle{Caractérisations par les familles{:}}}
Soit \(\mathscr{F} = (e_i)_{i \in \N}\) une famille de \(E\) et un endomorphisme de \(E\), alors \(f\) est entièrement caractérisée par l'image de la famille.

\subsection*{\subsecstyle{Endomorphismes remarquables{:}}}
Soit \(\varphi\) un endormorphisme de \(E\).\+
On considère ici des endomorphismes remarquables qui ont une interprétation géométrique trés intuitive comme:
\begin{itemize}
   \item Les homotéthies: Elles sont caractérisées par \(\varphi_k : u \longmapsto ku\)
   \item Les symétries: Elles sont caractérisées par \(\varphi \circ \varphi = Id_E\)
   \item Les projecteurs: Elles sont caractérisées par \(\varphi \circ \varphi = \varphi\)
\end{itemize}
Précisons le cas des projecteurs, en effet si on considère deux sous-espaces \(F, G\) supplémentaires dans \(E\), alors chaque élément \(u \in E\) admet une décomposition unique de la forme \(u = v_F + v_G\).\<

Alors le projecteur de direction \(F\) sur \(G\) est exactement l'application \(\varphi\) telle que:
\[
   \varphi(u) = v_G
\]
Graphiquement pour \(E = \R^2\), \(F, G\) deux droites vectorielles supplémentaires et \(\varphi\) le projecteur de direction \(F\) sur \(G\), on a:

\begin{center}
   \begin{tikzpicture}[]
      \draw[black!15] (-4,0) -- (4,0);
      \draw[black!15] (0, -2) -- (0,2);

      \draw[color=BrightBlue1, domain=-0.5:0.5, thick] plot (\x,{4*\x}) node[left] {$F$};
      \draw[color=BrightBlue1, domain=-4:4, thick] plot (\x,{0.5*\x}) node[below right] {$G$};

      \draw[-latex, color=BrightRed1, thick] (0, 0) -- (3/14, 6/7) node[left] {$v_F$};
      \draw[-latex, color=BrightRed1, thick] (0, 0) -- (9/7, 9/14) node[below right] {$v_G = \varphi$};

      \draw[color=BrightRed1, dashed] (9/7, 9/14) -- (1.5, 1.5);
      \draw[color=BrightRed1, dashed] (3/14, 6/7) -- (1.5, 1.5);

      \draw[-latex, color=DarkBlue1, thick] (0, 0) -- (1.5,1.5) node[above left] {$u$};

   \end{tikzpicture}
\end{center}

\subsection*{\subsecstyle{Applications multilinéaires{:}}}
On peut généraliser le concept de linéarité à celui de \textbf{multilinéarité}. On considère alors une application de la forme \(f : E^n \longrightarrow F\), alors on dit que \(f\) est multilinéaire si et seulement si elle est linéaire \textbf{en chacune des variables}, ie si:
\begin{align*}
   f(e_1 + \lambda u, e_2, \ldots, e_n) &= f(e_1, e_2, \ldots, e_n) + \lambda f(u, e_2, \ldots, e_n)  \\
   f(e_1, e_2  + \lambda u, \ldots, e_n) &= f(e_1, e_2, \ldots, e_n) + \lambda f(e_1, u, \ldots, e_n) \\
   & \vdotswithin{=} \\   
   f(e_1, e_2 , \ldots, e_n + \lambda u) &= f(e_1, e_2, \ldots, e_n) + \lambda f(e_1, e_2, \ldots, u)
\end{align*}
Une telle application est dite:
\begin{itemize}
   \item \textbf{Symétrique} si permuter deux variables préserve le résultat.\footnote[1]{Exemple: \(f(e_1, e_2) = f(e_2, e_1)\)}
   \item \textbf{Antisymétrique} si permuter deux variables change le signe du résultat.\footnote[2]{Exemple: \(f(e_1, e_2) = -f(e_2, e_1)\), en particulier, le signe du résultat aprés une permutation \(\sigma\) dépends alors de \textbf{la signature de la permutation} (la parité du nombre de permutations effectuées).}
   \item  \textbf{Altérnée} si elle s'annule à chaque fois qu'on l'évalue sur un k-uplet contenant deux vecteurs identiques.\footnote[3]{Exemple: \(f(e_1, e_1) = 0\)}
\end{itemize}
On peut alors définir la bilinéarité, la tri-linéarité ...

\chapter*{\chapterstyle{Algèbre II {:} Théorie de la dimension}}
Dans ce chapitre, on considérera un espace vectoriel \(E\) qui admet une famille génératrice \textbf{finie}. On dira alors que \(E\) est de \textbf{dimension finie}.\<

\subsection*{\subsecstyle{Théorème de la base incomplète{:}}}
Soit \(\mathcal{L}\) une famille libre et \(\mathcal{G}\) une famille génératrice de \(E\), le concept de \textbf{dimension} se définit gràce au \textbf{théorème de la base incomplète}:
\begin{itemize}
   \item On peut \textbf{compléter} \(\mathcal{L}\) en une base de \(E\) par ajouts de vecteurs de \(\mathcal{G}\).
   \item On peut \textbf{extraire} de \(\mathcal{G}\) une base de \(E\) (corollaire immédiat).
\end{itemize}
Ce théorème permet alors d'assurer \textbf{l'existence} d'une base d'une espace vectoriel de dimension finie.

\subsection*{\subsecstyle{Théorème de la dimension{:}}}
Il est évident que le cardinal d'une partie libre est toujours inférieur au cardinal d'une partie génératrice \footnote[1]{Aussi appellé \textbf{lemme de Steiniz}.}, de cette considération, on peut alors montrer \textbf{le théorème de la dimension}:
\begin{center}
   \textbox{Toutes les bases d'un espace vectoriel de dimensions fini ont \textbf{même cardinal}.}
\end{center}
Ce théorème permet alors d'assurer \textbf{l'unicité} du cardinal des bases d'une espace vectoriel de dimension finie.\<

\subsection*{\subsecstyle{Définition de la dimension{:}}}
Des deux théorèmes précédents, on a alors l'existence de bases d'une espace de dimension fini, et l'unicité de leur cardinal, on peut alors définir \textbf{la dimension d'une espace vectoriel} comme ce cardinal et on la note \(\text{dim}(E)\).\<
On peut alors en déduire la dimension des espaces vectoriels remarquables:
\begin{itemize}
   \item L'espace vectoriel classique \(\K^n\) est de dimension \(n\).
   \item L'espace vectoriel des polynômes de degré au plus \(n\) est de dimension \(n + 1\).
   \item Une droite vectorielle est de dimension \(1\).
   \item Un plan vectoriel est de dimension \(2\).
\end{itemize}

\subsection*{\subsecstyle{Espaces de dimension finie{:}}}
Considérons maintenant \(E\) un espace vectoriel de dimension finie \(n\) est \(\Fam = (e_i)_{i \in \N}\) une famille de \(n\) vecteurs de \(E\). Alors par déduction immédiate\footnote[2]{En effet par exemple sachant qu'elle est libre et possède \(n\) vecteurs, si ce n'était pas une base, on pourrait rajouter un/des vecteur pour obtenir une base, et donc dim\((E) \neq n \implies \bot\).} de la définition de dimension, on a:
\begin{center}
   \textbox{\(\Fam\) est libre \(\Longleftrightarrow\) \(\Fam\) est génératrice \(\Longleftrightarrow\) \(\Fam\) est une base}
\end{center}
Soient \(F, G\) deux sous-espaces de \(E\).\+
La dimension permet aussi de prouver des \textbf{égalités} d'espaces vectoriels, gràce à la propriété:
\[
   \begin{rcases}
      F \subset G \\
      \text{dim}(F) = \text{dim}(G)
   \end{rcases}\implies F = G
\]
Aussi, on peut alors \textbf{caractériser} le fait que \(F\) et \(G\) soient \textbf{supplémentaires} dans \(E\) par:
\[
   F \oplus G = E \Longleftrightarrow \begin{cases}
      F \cap G = \{0_E\}\\
      \text{dim}(F) + \text{dim}(G) = \text{dim}(E)
   \end{cases}
\]
Enfin on peut calculer \textbf{la dimension d'une somme}\footnote[1]{Le cas d'égalité est intéressant car il intervient pour des sous-espaces d'intersection nulle, et alors la dimension de la somme est simplement égale à la somme des dimensions.} avec la \textbf{formule de Grassmann}:
\[
   \mathbox{\text{dim}(F + G) = \text{dim}(F) + \text{dim}(G) - \text{dim}(F \cap G)}   
\]

\subsection*{\subsecstyle{Rang d'une famille{:}}}
Soit \(E\) un espace vectoriel de dimension finie et \(\Fam = (e_i)_{i \in \N}\) une famille de vecteurs de cet espace, alors on appelle \textbf{rang} de \(\Fam\) l'entier:
\[
   \mathbox{\text{rg}(\Fam) = \text{dim}(\text{Vect}(\Fam))}   
\]
Informellement, le rang est simplement \textbf{la dimension de l'espace engendré par la famille}.

\subsection*{\subsecstyle{Applications linéaires en dimension finie{:}}}
On appelle \textbf{rang d'une application linéaire la dimension de son image}, qu'on note \(\text{rg}(f)\), et on définit ainsi une application linéaire \textbf{de rang fini}.\<

Soit \(f \in \mathcal{L}(E, F)\) une telle application, alors on peut montrer un théorème fondamental appelé \textbf{théorème du rang}:
\[
   \mathbox{\text{dim}(E) = \text{rg}(f) + \text{dim}(\text{Ker}(f))}
\]
Ce thèorème permet alors de caractériser l'injectivité et la surjectivité d'une application linéaire \(\varphi\) par:
\begin{itemize}
   \item \(\varphi\) est injective si et seulement si \(\text{rg}(\varphi) = \text{dim}(E)\)
   \item \(\varphi\) est surjective si et seulement si \(\text{rg}(\varphi) = \text{dim}(F)\)
\end{itemize}

En particulier, si \(E\) et \(F\) sont deux espaces vectoriels \textbf{de même dimension} alors:
\begin{center}
   \textbox{\(\varphi\) est injective \(\Longleftrightarrow\) \(\varphi\) est surjective \(\Longleftrightarrow\) \(\varphi\) est bijective}
\end{center}
On caractérise ainsi \textbf{les isomorphismes en dimension finie}\footnote[2]{Toutes ces propriétés n'étant que de simples applications du théorème du rang.}.

\chapter*{\chapterstyle{Algèbre II {:} Espace des matrices}}
On appelle \textbf{matrice} à \(n\) lignes et \(p\) colonnes à coefficients dans \(\K\) toute application \(M : \inticc{1}{n} \times \inticc{1}{p} \longrightarrow \K\), en d'autres termes \textbf{tout famille de scalaires indexée par }\(\inticc{1}{n} \times \inticc{1}{p}\).\<

Informellement, il s'agit d'une généralisation du concept de suite sous forme de suite à \textbf{deux indices}, qu'on peut alors voir comme un tableau de nombres tel qu'en chaque position \((i, j)\), on ait un scalaire \(a_{ij} \in \K\).\< 

A l'instar des suites, on note \(M = (a_{ij})_{\substack{1 \leq i \leq n\\1 \leq j \leq p}}\) pour faire référence au terme général de \(M\) en position \((i,j)\).

\subsection*{\subsecstyle{Définition{:}}}
On appelle \textbf{espace des matrices} à \(n\) lignes et \(p\) colonnes à coefficient dans \(\K\), et on note \(\mathcal{M}_{n, p}(\K)\), l'espace de toutes les matrices comme définies ci-dessus.\<

On définit \textbf{la somme} de deux matrices comme la matrice dont le terme en position \((i, j)\) est la sommes des termes des deux matrices en cette position, informellement, on additionne simplement \textbf{composante par composante}.\+
On définit \textbf{la multiplication} d'une matrice par un scalaire comme la matrice dont tout les termes sont multipliés par ce scalaire.\<

Pour ces deux lois l'ensemble \((\mathcal{M}_{n, p}(\K), +, \cdot)\) est alors muni d'une structure \textbf{d'espace vectoriel} sur \(\K\) ayant pour vecteur nul \(0_E\) la matrice dont tout les coefficients sont nuls..

\subsection*{\subsecstyle{Produit matriciel{:}}}
On définit aussi une multiplication entre matrices compatibles au sens ci-dessous\footnote[1]{Simplement, il faut que le nombre de colonnes de la première soit égal au nombre de lignes de la seconde.}.\<

Formellement, soit \(A = (a_{i, j}) \in \mathcal{M}_{n, m}(\K)\) et \(B = (b_{k, j}) \in \mathcal{M}_{m, p}(\K)\).\+
Alors on définit la matrice \(C := AB = (c_{i,j})\) comme étant la matrice telle que:
\[
   \mathbox{c_{i, j} = \sum_{k=1}^{n}a_{i, k}b_{k, j}}   
\]
En particulier, on appelle ce produit \textbf{un produit ligne par colonne} qui se comprends visuellement par:\+
\begin{center}
   \begin{tikzpicture}[    
      every left delimiter/.style={xshift=.55em},
      every right delimiter/.style={xshift=-.55em}
   ]
      \matrix[
         matrix of math nodes, ampersand replacement=\&,
         left delimiter=(, right delimiter=),  outer sep = 0pt,inner sep=4.5pt
      ](A){
         a_{1, 1} \& a_{1, 2}\\ 
         a_{2, 1} \& a_{2, 2}\\ 
         a_{3, 1} \& a_{3, 2}\\  
      };
      \draw[color=BrightBlue1, line width=0.5mm] (A-3-1.south west) rectangle (A-3-2.north east);
      \draw node at (1.5, 0) {\(\times\)};

      \matrix[
         matrix of math nodes, ampersand replacement=\&,
         left delimiter=(, right delimiter=),  outer sep = 0pt,inner sep=4.5pt
      ] (B) at (3.4, 0) {
         b_{1, 1} \& b_{1, 2} \& b_{1, 3}\\ 
         b_{2, 1} \& b_{2, 2} \& b_{2, 3}\\ 
      };
      \draw[color=BrightBlue1, line width=0.5mm] (B-1-1.north west) rectangle (B-2-1.south east);

      \draw node at (5.25, 0) {\(=\)};
   
      \matrix[
         matrix of math nodes, ampersand replacement=\&,
         left delimiter=(, right delimiter=),  outer sep = 0pt,inner sep=4.5pt
      ] (C) at (7.1, 0) {
         c_{1, 1} \& c_{1, 2} \& c_{1, 3}\\ 
         c_{2, 1} \& c_{2, 2} \& c_{2, 3}\\ 
         c_{3, 1} \& c_{3, 2} \& c_{3, 3}\\
      };
      \draw[color=BrightRed1, line width=0.5mm] (C-3-1.north west) rectangle (C-3-1.south east);
   \end{tikzpicture}   
\end{center}
Le coefficient à \color{BrightRed1} la troisième ligne, première colonne \color{black} est obtenu en multipliant \color{BrightBlue1} la troisième ligne par la première colonne.\color{black}\<

On appelle \textbf{matrice identité} la matrice neutre pour cette opération, c'est la matrice composée de \(1\) sur toute la diagonale principale, et de \(0\) partout ailleurs.\<

En considérant cette définition, il est évident que le produit matrice n'est en général \textbf{pas commutatif}. Par ailleurs \textbf{le produit de deux matrices non nulles peut être nul}.
\pagebreak

Alors sous réserve de compatibilités des tailles matrices, on a les propriétés suivantes:
\begin{multicols}{2}
   \begin{itemize}
      \item Associativité
      \item La matrice identité est neutre
      \item Distributivité
      \item La matrice nulle est absorbante
   \end{itemize}
\end{multicols}   

Considérons maintenant le cas des matrices \textbf{carrées}, le produit avec elle même est toujours défini, donc on peut définir une \textbf{puissance positive de matrice}, qu'on note \(A^n\) qui vérifie alors:

\begin{multicols}{2}
   \begin{itemize}
      \item \((A^p)^q = A^{pq}\)
      \item \(A^pA^q = A^{p +q}\)
   \end{itemize}
\end{multicols}    

\subsection*{\subsecstyle{Transposition{:}}}
Soit \(M = (x_{i,j})\in \mathcal{M}_{n, m}(\K)\), on définit maintenant \textbf{une application linéaire involutive} fondamentale de l'espace des matrices, il s'agit de \textbf{l'opération de transposition} notée \(M^\mathsf{T}\) et définie par:
\[
   \mathbox{M^\mathsf{T} = (x_{j,i})}   
\]
Intuitivement, cette application transforme \textbf{chaque ligne en colonne et invérsement}. Par exemple:
\[
   \begin{pmatrix}
      a & b & c\\
      d & e & f
   \end{pmatrix}^\top = 
   \begin{pmatrix}
      a & d\\
      b & e\\
      c & f
   \end{pmatrix}
\]
Il faut aussi noter son comportement par rapport au produit matriciel de deux matrices \(A, B\), on a:
\[
   \mathbox{
      (AB)^\top = B^\top A^\top
   }
\]

\subsection*{\subsecstyle{Trace{:}}}
On définit aussi une autre application linéaire sur l'espace des matrices appellée \textbf{trace d'une matrice}, et qui est définie comme \textbf{la somme des éléments diagonaux}. Formellement:
\[
   \text{tr}(A) := \sum_{i=1}^{n} a_{i, i}   
\]
Elle est donc linéaire mais on a aussi:
\[
   \text{tr}(AB) = \text{tr}(BA)   
\]

\subsection*{\subsecstyle{Matrice d'une famille de vecteurs{:}}}
Soit \(\mathscr{F} = (e_i)_{i\in\N}\) une famille de vecteurs, alors on peut définir \textbf{la matrice de la famille} dans la base canonique par:
\[
   \text{Mat}_{\mathscr{C}}(\mathscr{F}) = ([e_1]_{\mathscr{C}}, [e_2]_{\mathscr{C}}, \ldots, [e_n]_{\mathscr{C}})  
\]
En d'autres termes, c'est simplement \textbf{la matrice constituée des coordonnées de vecteurs de la famille} rangées en colonnes.\<

Alors en particulier, on montre aisément, que \(Vect()\)

A FINIR

\subsection*{\subsecstyle{Matrice d'une application linéaire{:}}}
Une des utilités fondamentales de l'espace des matrices et son lien puissant avec \textbf{les applications linéaires}.\<

Soit \(E, F\) deux espaces vectoriels de bases \(\mathscr{B} = (e_i)_{i \in \N}\) et \(\mathscr{C}=(f_i)_{i \in \N}\).\+
On considère une application linéaire \(f : E \longrightarrow F\), alors on peut associer à l'application \(f\) \textbf{une unique matrice} dans les bases \(\mathscr{B, C}\), qu'on note alors Mat\(_{\mathscr{B, C}}(f)\) qu'on construit comme suit:
\[
   \mathbox{\text{Mat}_{\mathscr{B, C}}(f) = ([f(e_1)]_{\mathscr{C}}, [f(e_2)]_{\mathscr{C}}, \ldots, [f(e_n)]_{\mathscr{C}})}   
\]
Ainsi, il s'agit de la matrice dont la i-ème colonne comporte les coordonnées dans la base \(\mathscr{C}\) de l'image de i-ème vecteur de la base \(\mathscr{B}\).
Informellement:
\begin{center}
   \textit{
      La matrice d'une application linéaire est constituée des coordonées de l'image des vecteurs de la base de départ (en colonnes) dans la base d'arrivée.
   }
\end{center}

Exemple: Considérons un endomorphisme de \(\R^3\) et sa base canonique notée \((e_1, e_2, e_3)\), telle que \(f(x, y, z) = (2x+1, 3x, z-1)\). Alors on calcule l'image des vecteurs de la base de départ, ie:
\begin{multicols}{3}
\begin{itemize}
   \item f(1, 0, 0) = (3, 3, -1)
   \item f(0, 1, 0) = (0, 0, -1)
   \item f(0, 0, 1) = (0, 0, -2)
\end{itemize}
\end{multicols}
Puis on calcule les coordonées de ces vecteurs dans la base d'arrivée, et on les range en colonne dans une matrice et on obtient:\+

\begin{center}
   \begin{tikzpicture}[      
      every left delimiter/.style={xshift=.55em},
      every right delimiter/.style={xshift=-.55em}
   ]
   \draw node[color = BrightBlue1] at (-1.05, 1.75) {$f(e_1)$};
   \draw[-latex, dashed, thick, color = BrightBlue1!50] (-1.05, 1.5) -- (-1.05, 1.1);

   \draw node[color = BrightBlue1] at (0, 1.75) {$f(e_2)$};
   \draw[-latex, dashed, thick, color = BrightBlue1!50] (0, 1.5) -- (0, 1.1);

   \draw node[color = BrightBlue1] at (1.05, 1.75) {$f(e_3)$};
   \draw[-latex, dashed, thick, color = BrightBlue1!50] (1.05, 1.5) -- (1.05, 1.1);

   \draw node[color = BrightBlue1] at (2.3, 0.85) {$f_1$};
   \draw[-latex, dashed, thick, color = BrightBlue1!50] (1.7, 0.85) -- (2.1, 0.85);

   \draw node[color = BrightBlue1] at (2.3, 0) {$f_2$};
   \draw[-latex, dashed, thick, color = BrightBlue1!50] (1.7, 0) -- (2.1, 0);

   \draw node[color = BrightBlue1] at (2.3, -0.85) {$f_3$};
   \draw[-latex, dashed, thick, color = BrightBlue1!50] (1.7, -0.85) -- (2.1, -0.85);

   \matrix[
      matrix of math nodes, ampersand replacement=\&,
      left delimiter=(, right delimiter=), inner ysep=4pt, column sep=0.8em, row sep = 0.8em,
      nodes={
         inner sep=0.5em,outer sep=0,text width=1.2em,align=center
      }
   ]{
      3 \& 0 \& 0\\ 
      3 \& 0 \& 0\\ 
      -1 \& -1 \& -2\\
   };
   \end{tikzpicture}
\end{center}
Cette construction implique qu'il suffit donc, pour une base donnée, de \textbf{calculer les images des vecteurs de cette base}, puis leur coordonées pour décrire la transformation.

Soit \(E\) et \(F\) sont de dimensions respectives \(n\) et \(p\), alors on peut montrer le théorème le plus important de cette partie, qui exhibe \textbf{un isomorphisme fondamental}:
\[
   \mathbox{
      \text{Mat}_{\mathscr{B, C}} : \mathcal{L}(E, F) \xlongrightarrow{\sim} \mathcal{M}_{n, p}(\K)  
   }
\]
On peut alors considèrer, que du point de vue de l'algèbre linéaire:
\begin{center}
   \textit{L'espace des application linéaires et l'espace des matrices sont les mêmes\footnote[1]{En dimension finie}.}
\end{center}
En particulier:
\[
   f = g \Longleftrightarrow \text{Mat}_{\mathscr{B}}(f) = \text{Mat}_{\mathscr{B}}(g)
\]

Dans la suite, pour plus de lisibilité, on considèrera un endomorphisme \(f\) de \(\R^2\) de matrice \(M\) dans la base canonique \(\mathscr{C}\).
Alors, appliquer une transformation linéaire à un vecteur \(u \in E\) revient à \textbf{multiplier ce vecteur par la matrice associée}\footnote[2]{Il suffit de partir du produit matriciel et d'utiliser la définition de \(M\) comme étant la matrice de \(f\) et les règles de calculs sur les vecteurs colonnes.}, ie on a:
\[
   \mathbox{
      [f(u)]_{\mathscr{C}} = M \times [u]_{\mathscr{C}}
   }
\]
Exemple: Soit \(f\) de matrice 
\(M = \begin{pmatrix}
   1 & 2\\
   3 & 4\\
\end{pmatrix}\) 
et \(u = (5, 3)\), calculer les coordonées de \(f(u)\) revient à calculer:
\[
   \begin{pmatrix}
      1 & 2\\
      3 & 4\\
   \end{pmatrix}
   \begin{pmatrix}
      5 \\
      3 \\
   \end{pmatrix} =
   \begin{pmatrix}
      11 \\
      27 \\
   \end{pmatrix} 
\]
Enfin l'isomorphisme entre ces deux espaces nous permet de définir \textbf{le noyau, l'image, et le rang d'une matrice}, comme simplement étant le noyau, l'image, ou le rang de l'application linéaire associée à cette matrice.

\subsection*{\subsecstyle{Matrices inversibles{:}}}
Soit \(M \in \mathcal{M}_n(\K)\) une matrice carrée, alors on dit que \(M\) est \textbf{inversible} si il existe un inverse pour la loi de multiplication des matrices, ie si il existe une matrice \(A^{-1}\) telle que:
\[
   \mathbox{AA^{-1}=Id_n}   
\]
En particulier, si on considère l'application linéaire associée à \(M\), alors:
\begin{center}
   \(M\) inversible \(\Longleftrightarrow\) \(f\) est inversible
\end{center}\<

On appelle l'ensemble des matrices carrées inversibles de taille \(n\) \textbf{le groupe linéaire} d'ordre \(n\), qu'on note GL\(_n(\K)\), c'est un groupe pour la multiplication matricielle qui est stable par produit, ie:
\[
   A, B \in \text{GL}_n(\K) \Longleftrightarrow AB \in \text{GL}_n(\K)
\]
Et en particulier, on montre\footnote[2]{Intuitivement, c'est analogue à l'inverse d'une composée car les matrices "sont" des applications au sens de l'algèbre linéaire.} que \(AB^{-1} = B^{-1}A^{-1}\). Enfin on peut définir les puissances(entières) négatives d'une matrice inversible de manière analogue à la définition des puissances positives.\<

Intuitivement:
\begin{center}
   \textit{Le groupe \(\text{GL}_n(\K)\) est donc le groupe \textbf{des automorphismes} de \(E\), c'est à dire le groupe \textbf{des transformations} (linéaires) de \(E\) qui sont \textbf{réversibles}.}
\end{center}
Visuellement, on comprends directement l'idée d'automorphisme d'un espace, son lien avec la bijectivité, et la matrice inverse est alors la représentation algèbrique de la transformation inverse (ici avec une rotation de \(\R^2\) de \(45\) degrés):
\begin{center}
   \begin{tikzpicture}
      \draw[step=0.45cm, color=black!20, dashed] (-1.75,-1.75) grid (1.75,1.75);
      \draw[-latex, color=BrightBlue1, thick] (0,0) -- (0,1.5) node[below left] {$y$};
      \draw[-latex, color=BrightBlue1, thick] (0,0) -- (1.5,0) node[below left] {$x$};
      \draw[-latex, color=BrightRed1, thick] (0,0) -- (1,1.2) node[below right] {$u$};
   
      \draw[-latex] (2,0) -- (3,0);
      \draw[] (2,-0.05) -- (2,0.05);
      \draw[] node[] at (2.5,0.3){$\varphi$};
   
      \draw[step=0.45cm, color=black!20, dashed, shift={(5.5 cm, 0)}, rotate=45] (-1.75,-1.75) grid (1.75,1.75);
      \draw[-latex, color=BrightBlue1, thick, shift={(5.5 cm, 0)}, rotate=45] (0,0) -- (0,1.5) node[below left] {$\varphi(y)$};
      \draw[-latex, color=BrightBlue1, thick, shift={(5.5 cm, 0)}, rotate=45] (0,0) -- (1.5,0) node[below right] {$\varphi(x)$};
      \draw[-latex, color=BrightRed1, thick, shift={(5.5 cm, 0)}, rotate=45] (0,0) -- (1,1.2) node[below right] {$\varphi(u)$};
   
      \draw[-latex] (8,0) -- (9,0);
      \draw[] (8,-0.05) -- (8,0.05);
      \draw[] node[] at (8.65,0.375){$\varphi^{-1}$};
   
      \draw[step=0.45cm, color=black!20, dashed, shift={(11 cm, 0)}] (-1.75,-1.75) grid (1.75,1.75);
      \draw[-latex, color=BrightBlue1, thick, shift={(11 cm, 0)}] (0,0) -- (0,1.5) node[below left] {$y$};
      \draw[-latex, color=BrightBlue1, thick, shift={(11 cm, 0)}] (0,0) -- (1.5,0) node[below left] {$x$};
      \draw[-latex, color=BrightRed1, thick, shift={(11 cm, 0)}] (0,0) -- (1,1.2) node[below right] {$u$};
   \end{tikzpicture}   
\end{center}
Enfin, on peut montrer que \textbf{la transposition} est compatible avec l'inverse, ie on a:
\[
   \mathbox{(M^{-1})^{\top} = (M^{\top})^{-1}}   
\]
Montrer qu'une matrice est inversible ou non peut se réaliser de plusieurs manières, on peut utliser le lien avec les applications linéaires (montrer que \(f\) est bijective) qui revient à résoudre un systême ou alors utiliser une méthode d'échelonnement décrite dans le chapitre associé.\<

On définira aussi une application fondamentale nommée \textbf{déterminant} qui sera définie dans la prochaine partie et qui caractérise exactement toutes les matrices inversibles.

\subsection*{\subsecstyle{Matrices semblables{:}}}
On considère un espace vectoriel \(E\) et deux bases \(\Fam=(e_1, \ldots, e_n), \Fam'=({e}_1', \ldots, {e}_n')\) de cet espace. Dans la suite, on considérera \(n = 3\) pour plus de lisibilité.\<

On appelle \textbf{matrice de passage} de \(\Fam\) à \(\Fam'\) la matrice \(P\) qui vérifie:
\begin{center}
   \textbox{\textit{Ses colonnes sont les coordonées de \textbf{la nouvelle base} dans \textbf{l'ancienne base}}.}
\end{center}
Visuellement, on a:

\begin{center}
   \begin{tikzpicture}[      
      every left delimiter/.style={xshift=.55em},
      every right delimiter/.style={xshift=-.55em}
   ]
   \draw node[color = BrightBlue1] at (-1.05, 1.75) {$e_1'$};
   \draw[-latex, dashed, thick, color = BrightBlue1!50] (-1.05, 1.5) -- (-1.05, 1.1);

   \draw node[color = BrightBlue1] at (0, 1.75) {$e_2'$};
   \draw[-latex, dashed, thick, color = BrightBlue1!50] (0, 1.5) -- (0, 1.1);

   \draw node[color = BrightBlue1] at (1.05, 1.75) {$e_3'$};
   \draw[-latex, dashed, thick, color = BrightBlue1!50] (1.05, 1.5) -- (1.05, 1.1);

   \draw node[color = BrightBlue1] at (2.3, 0.85) {$e_1$};
   \draw[-latex, dashed, thick, color = BrightBlue1!50] (1.7, 0.85) -- (2.1, 0.85);

   \draw node[color = BrightBlue1] at (2.3, 0) {$e_2$};
   \draw[-latex, dashed, thick, color = BrightBlue1!50] (1.7, 0) -- (2.1, 0);

   \draw node[color = BrightBlue1] at (2.3, -0.85) {$e_3$};
   \draw[-latex, dashed, thick, color = BrightBlue1!50] (1.7, -0.85) -- (2.1, -0.85);

   \matrix[
      matrix of math nodes, ampersand replacement=\&,
      left delimiter=(, right delimiter=), inner ysep=4pt, column sep=0.8em, row sep = 0.8em,
      nodes={
         inner sep=0.5em,outer sep=0,text width=1.2em,align=center
      }
   ]{
      * \& * \& *\\ 
      * \& * \& *\\ 
      * \& * \& *\\
   };
   \end{tikzpicture}
\end{center}

Une matrice de passage est alors \textbf{nécessairement inversible}, et si on note \(U = [u]_{\Fam}\) et \(U' = [u]_{\Fam'}\)on a alors le théorème fondamental suivant:
\[
   \mathbox{U' = P^{-1}U}
\]
En d'autres termes:
\begin{center}
   \textit{Pour obtenir les coordonées d'un vecteur dans la nouvelle base, il suffit de le multiplier par l'inverse de la matrice de passage.}
\end{center}
Considérons maintenant un endomorphisme \(f\) de matrices \(M = \text{Mat}(\Fam, f), M' = \text{Mat}(\Fam', f)\) dans les deux bases respectivement, alors on a:
\[
   \mathbox{M' = P^{-1}MP
\]
Enfin on dira alors que deux matrices \(A, B \in \mathcal{M}_n(\K)\) sont \textbf{semblables} si il existe une matrice de passage telle que la relation ci-dessus soit vérifiée.
Ce qui signifie exactement que:
\begin{center}
   \textit{Deux matrices semblables représentent \textbf{la même transformation géométrique} dans des bases différentes.}
\end{center}
En particulier on appelle alors \textbf{invariants de similitude} les propriétés qui \textbf{ne dépendent pas du choix de la base}, on peut alors montrer que:
\begin{center}
   \textit{Le rang et la trace d'une matrice (et donc d'une application linéaire) sont des invariants de similitude.}
\end{center}

\pagebreak
\subsection*{\subsecstyle{Matrices remarquables{:}}}
Dans ce chapitre nous allons présenter brièvement différentes matrices remarquables:
\begin{figure}[h]
   \color{DarkBlue1}
   \centering
   \begin{subfigure}{.3\textwidth}
      \centering
      \begin{tikzpicture}[line cap=round]
         \matrix[
            matrix of math nodes, ampersand replacement=\&,
            left delimiter=(, right delimiter=)
         ]{
            a \& 0 \& 0\\ 
            0 \& d \& 0\\ 
            0 \& 0 \& f\\
         };
      \end{tikzpicture}
      \caption*{\color{DarkBlue1}
      Matrice diagonale}
   \end{subfigure}\quad
   \begin{subfigure}{.3\textwidth}
      \centering
      \begin{tikzpicture}[line cap=round]
         \matrix[
            matrix of math nodes, ampersand replacement=\&,
            left delimiter=(, right delimiter=)
         ]{
            a \& b \& c\\ 
            0 \& d \& e\\ 
            0 \& 0 \& f\\
         };
      \end{tikzpicture}
      \caption*{\color{DarkBlue1}
      Matrice triangulaire supérieure}
   \end{subfigure}\quad
   \begin{subfigure}{.3\textwidth}
      \centering
      \begin{tikzpicture}[line cap=round]
         \matrix[
            matrix of math nodes, ampersand replacement=\&,
            left delimiter=(, right delimiter=)
         ]{
            a \& 0 \& 0\\ 
            b \& d \& 0\\ 
            c \& e \& f\\
         };
      \end{tikzpicture}
      \caption*{\color{DarkBlue1}
      Matrice triangulaire inférieure}
   \end{subfigure}
\end{figure}

Puis on a deux types de matrices qui sont liées à la symétrie des coefficients et à \textbf{l'opération de transposition}\footnote[1]{En effet une matrice symétrique est égale à sa transposée et une matrice antisymétrique à \textbf{l'opposée} de sa transposée}:
\begin{figure}[h]
   \color{DarkBlue1}
   \centering
   \begin{subfigure}{.3\textwidth}
      \centering
      \begin{tikzpicture}[line cap=round]
         \matrix[
            matrix of math nodes, ampersand replacement=\&,
            left delimiter=(, right delimiter=)
         ]{
            a \& d \& f\\ 
            d \& b \& e\\ 
            f \& e \& c\\
         };
      \end{tikzpicture}
      \caption*{\color{DarkBlue1}Matrice symétrique}
   \end{subfigure}\quad
   \begin{subfigure}{.3\textwidth}
      \centering
      \begin{tikzpicture}[line cap=round]
         \matrix[
            matrix of math nodes, ampersand replacement=\&,
            left delimiter=(, right delimiter=)
         ]{
            0 \& -a \& -b\\ 
            a \& 0 \& -c\\ 
            b \& c \& 0\\
         };
      \end{tikzpicture}
      \caption*{\color{DarkBlue1}Matrice antisymétrique}
   \end{subfigure}
\end{figure}

Enfin on a \textbf{les matrices élémentaires} qui sont obtenus par \textbf{opérations élémentaires}\footnote[1]{Voir chapitre sur l'échelonnement.} sur la matrice identité:
\begin{figure}[h]
   \color{DarkBlue1}
   \centering
   \begin{subfigure}{.3\textwidth}
      \centering
      \begin{tikzpicture}[line cap=round]
         \matrix[
            matrix of math nodes, ampersand replacement=\&,
            left delimiter=(, right delimiter=)
         ]{
            1 \& 0 \& 0\\ 
            0 \& 1 \& 0\\ 
            0 \& 0 \& \lambda\\
         };
      \end{tikzpicture}
      \caption*{\color{DarkBlue1}Matrice de dilatation}
   \end{subfigure}\quad
   \begin{subfigure}{.3\textwidth}
      \centering
      \begin{tikzpicture}[line cap=round]
         \matrix[
            matrix of math nodes, ampersand replacement=\&,
            left delimiter=(, right delimiter=)
         ]{
            1 \& 0 \& 0\\ 
            0 \& 1 \& 0\\ 
            0 \& \lambda\& 1\\
         };
      \end{tikzpicture}
      \caption*{\color{DarkBlue1}Matrice de transvection}
   \end{subfigure}\quad
   \begin{subfigure}{.3\textwidth}
      \centering
      \begin{tikzpicture}[line cap=round]
         \matrix[
            matrix of math nodes, ampersand replacement=\&,
            left delimiter=(, right delimiter=)
         ]{
            1 \& 0 \& 0\\ 
            0 \& 0 \& 1\\ 
            0 \& 1\& 0\\
         };
      \end{tikzpicture}
      \caption*{\color{DarkBlue1}Matrice de permutation}
   \end{subfigure}
\end{figure}

\subsection*{\subsecstyle{Echelonnements \& Calculs{:}}}
La théorie permettant de lier, matrices, applications linéaires et familles de vecteurs, l'étude des espaces vectoriels engendrés par les colonnes ou les lignes d'une matrice revêt alors une importantce capitale qui est l'objet de ce chapitre.\<

On appelle \textbf{opérations élémentaires} sur une matrice \(M\), l'une de ces 3 opérations:
\begin{itemize}
   \item Multiplier une colonne par un scalaire
   \item Ajouter une colonnes à une autre
   \item Echanger deux colonnes
\end{itemize}
Il est alors possible de montrer que ces trois opérations préservent \textbf{le sous-espace engendré}\footnote[3]{L'échange ne change évidement rien, et le sous espace engendré ne change pas pas multiplication par un scalaire ou ajout d'un vecteur venant de la famille en question} par les colonnes de la matrice, donc en particulier \textbf{le rang, le noyau et l'image} sont préservés.\<

Enfin, on appellera \textbf{matrice équivalente} à \(M\) et on notera \(M' \sim M\) la matrice \(M\) à laquelle on a appliqué une ou plusieurs opérations élémentaires.
\pagebreak

On dira qu'une matrice \(M\) est \textbf{échelonée} si la matrice obtenue aprés application d'opérations élémentaires est de la forme générale:
\begin{center}
   \begin{tikzpicture}[    
      every left delimiter/.style={xshift=.55em},
      every right delimiter/.style={xshift=-.55em}
   ]
      \matrix[
         matrix of math nodes, ampersand replacement=\&,
         left delimiter=(, right delimiter=),  outer sep = 0pt,inner sep=4.5pt
      ] (A) {
         {\alpha_1} \& {0} \& {0} \& {0}\& {0}\\ 
         {*} \& {\alpha_2} \& {0} \& {0}\& {0}\\ 
         {*} \& {*} \& {0} \& {0}\& {0}\\ 
         {*} \& {*} \& {\alpha_3} \& {0}\& {0}\\
      };
      \draw[color=BrightBlue1, line width=0.5mm] (A-1-1.north west) -- (A-1-1.north east) -- (A-1-1.south east) -- (-0.1, 0.57)-- (-0.1, -0.64) -- (A-4-3.north east)-- (A-4-3.south east);

   \end{tikzpicture}   
\end{center}
En d'autres termes, on \textbf{creuse} la matrice pour faire apparaître des zéros sur la partie triangulaire supérieure
Exemple: Soit \(M = \begin{pmatrix}
   1 & 2 & 3\\
   4 & 5 & 6\\
   7 & 8 & 9\\
\end{pmatrix}\), alors on a:
\begin{flalign*}
   M =  \begin{pmatrix}  1 & 2 & 3\\ 4 & 5 & 6\\ 7 & 8 & 9\end{pmatrix} \sim\footnote[1]{Par exemple ici on a effectué \(C_2 \leftarrow C_2 - 2C_1\) et \(C_3 \leftarrow C_3 - 3C_1\)} \begin{pmatrix}1 & 0 & 0\\ 4 & -3 & -6\\ 7 & -6 & -12\\\end{pmatrix} \sim \begin{pmatrix}1 & 0 & 0\\ 4 & -3 & 0\\ 7 & -6 & 0\\\end{pmatrix}\shorteqnote{\color{BrightRed1}(\(\star\))}
\end{flalign*}
Il existe aussi une variante appellé \textbf{échelonnement avec mémorisation}, pour cette variante, on nomme chaque vecteur-colonne de la matrice et on reporte toutes les transformations réalisées sur ces vecteurs.
Exemple: Pour la même matrice, alors on a:
\begin{flalign*}
   M =
   \begin{blockarray}{ccc}
      C_1 & C_2 & C_3 \\
      \begin{block}{(ccc)}
         1 & 2 & 3\\
         4 & 5 & 6\\
         7 & 8 & 9\\
      \end{block}
   \end{blockarray} \sim
   \begin{blockarray}{ccc}
      C_1 & C_2-2C_1 & C_3-3C_3 \\
      \begin{block}{(ccc)}
         1 & 0 & 0\\
         4 & -3 & -6\\
         7 & -6 & -12\\
      \end{block}
   \end{blockarray}\sim\begin{blockarray}{ccc}
                          C_1 & C_2-2C_1 & C_3-3C_1 - 2C_2\\
                          \begin{block}{(ccc)}
         1 & 0 & 0\\
         4 & -3 & 0\\
         7 & -6 & 0\\
      \end{block}
                       \end{blockarray}\shorteqnote{\color{BrightRed1}(\(\star\star\))}
\end{flalign*}
Présentons maintenant les différentes informations que nous pouvons tirer de ces échellonements:


\chapter*{\chapterstyle{Algèbre II {:} Déterminant}}

\subsection*{\subsecstyle{Propriétés{:}}}

\subsection*{\subsecstyle{Cofacteurs \& Comatrices{:}}}

\subsection*{\subsecstyle{Formules de Cramer{:}}}

